\documentclass[11pt,a4paper]{article}
\usepackage{fontspec,amsmath,amssymb,amsthm,microtype}
\setmainfont{lmroman10-regular.otf}[BoldFont=lmroman10-bold.otf,ItalicFont=lmroman10-italic.otf,BoldItalicFont=lmroman10-bolditalic.otf]
\usepackage{xcolor}
\usepackage[most]{tcolorbox}
\usepackage{mathrsfs}
\providecommand{\lt}{<}
\providecommand{\gt}{>}
\setmonofont{DejaVuSansMono.ttf}[BoldFont=DejaVuSansMono-Bold.ttf,ItalicFont=DejaVuSansMono-Oblique.ttf,BoldItalicFont=DejaVuSansMono-BoldOblique.ttf,Scale=0.86]
\newfontfamily\leanfallback{STIXTwoMath-Regular.otf}[Scale=0.86]
\definecolor{leanrule}{gray}{0.62}
\definecolor{leanback}{gray}{0.962}
\definecolor{leantitle}{gray}{0.22}
\definecolor{leanmuted}{gray}{0.38}
\newtcolorbox{leanbox}[1]{enhanced,breakable,colback=leanback,colframe=leanrule,
  boxrule=0.4pt,arc=0.6mm,left=2.4mm,right=2.4mm,top=1.8mm,bottom=1.8mm,
  before skip=0.8em,after skip=1em,fonttitle=\small\bfseries,coltitle=leantitle,
  colbacktitle=leanback,titlerule=0.3pt,title={#1}}
\newlength{\leancharwidth}
\newenvironment{leancode}{\par\smallskip\begingroup\ttfamily\footnotesize\raggedright
  \settowidth{\leancharwidth}{M}\setlength{\parskip}{0pt}\setlength{\parindent}{0pt}
  \frenchspacing}{\par\endgroup\smallskip}
\newcommand{\leanline}[2]{\hangindent=\dimexpr#1\leancharwidth+4\leancharwidth\relax
  \hangafter=1\noindent\hspace*{#1\leancharwidth}#2\par}
\newcommand{\leansource}[1]{\par{\footnotesize\color{leanmuted}#1}\par}
\newcommand{\leanlabel}[1]{\par\smallskip\noindent{\small\bfseries #1}\par}
\usepackage[margin=25mm]{geometry}
\usepackage[unicode,hidelinks]{hyperref}
\hypersetup{pdfauthor={LangevinLean Project}}
\setlength{\emergencystretch}{3em}
\setlength{\parskip}{0.4em}
\setlength{\parindent}{0pt}
\allowdisplaybreaks[2]
\raggedbottom
\title{\texorpdfstring{Continuity of actual stationary observable evaluations}{Continuity of actual stationary observable evaluations}}
\author{LangevinLean Project}
\date{}
\begin{document}\maketitle
Draft mathematical exposition; these notes may contain mistakes.


If stationary random states with common law \(\mu\) converge in measure, their evaluation maps converge strongly on every \(L^2(\mu)\) observable. The proof first treats bounded continuous functions and then uses their density and the common evaluation isometry.


\section{\texorpdfstring{The actual evaluation and its square norm}{The actual evaluation and its square norm}}


\textbf{Claim.} A state \(X\) with law \(\mu\) gives an evaluation isometry \(f\mapsto f(X)\). For bounded continuous \(\phi\),

\[\|\phi(X)-\phi(Y)\|_2^2=\int\|\phi(X(\omega))-\phi(Y(\omega))\|^2\,dQ.\]

\textbf{Proof.} Evaluation is pullback by a measure-preserving state map. Compare the exact almost-everywhere representatives and use the integral square-norm formula. The auxiliary implication from vanishing mean norm to convergence in measure follows from the integral Markov bound.

\begin{leanbox}{Lean: \texttt{tendstoInMeasure\_\allowbreak{}of\_\allowbreak{}integral\_\allowbreak{}norm\_\allowbreak{}sub}, \texttt{evaluation}, \texttt{lp\_\allowbreak{}norm\_\allowbreak{}sq\_\allowbreak{}eq\_\allowbreak{}integral} and 2 more}
\leansource{\texttt{tendstoInMeasure\_\allowbreak{}of\_\allowbreak{}integral\_\allowbreak{}norm\_\allowbreak{}sub} \quad WassersteinMixing/\allowbreak{}SemigroupContinuity.lean, lines 30–43}
\begin{leancode}
\leanline{0}{theorem\ tendstoInMeasure\_of\_integral\_norm\_sub\ \{I\ :\ Type*\}\ \{l\ :\ Filter\ I\}}
\leanline{4}{(X\ :\ I\ →\ Ω\ →\ E)\ (Y\ :\ Ω\ →\ E)}
\leanline{4}{(hint\ :\ ∀\ i,\ Integrable\ (fun\ ω\ ↦\ X\ i\ ω\ -\ Y\ ω)\ Q)}
\leanline{4}{(hlim\ :\ Tendsto\ (fun\ i\ ↦\ ∫\ ω,\ ‖X\ i\ ω\ -\ Y\ ω‖\ ∂Q)\ l\ ({\leanfallback 𝓝}\ 0))\ :}
\leanline{4}{TendstoInMeasure\ Q\ X\ l\ Y}
\end{leancode}
\leansource{\texttt{evaluation} \quad WassersteinMixing/\allowbreak{}SemigroupContinuity.lean, lines 46–49}
\begin{leancode}
\leanline{0}{def\ evaluation\ (X\ :\ Lp\ E\ 2\ Q)\ (hX\ :\ Q.map\ X\ =\ μ)\ :\ Lp\ F\ 2\ μ\ →ₗᵢ[{\leanfallback 𝕜}]\ Lp\ F\ 2\ Q\ :=}
\leanline{2}{Lp.compMeasurePreservingₗᵢ\ {\leanfallback 𝕜}\ X\ ⟨(Lp.stronglyMeasurable\ X).measurable,\ hX⟩}
\leanline{0}{\strut}
\leanline{0}{include\ {\leanfallback 𝕜}\ in}
\end{leancode}
\leansource{\texttt{lp\_\allowbreak{}norm\_\allowbreak{}sq\_\allowbreak{}eq\_\allowbreak{}integral} \quad WassersteinMixing/\allowbreak{}SemigroupContinuity.lean, lines 52–58}
\begin{leancode}
\leanline{0}{theorem\ lp\_norm\_sq\_eq\_integral\ (Z\ :\ Lp\ F\ 2\ Q)\ :}
\leanline{4}{‖Z‖\ \textasciicircum{}\ 2\ =\ ∫\ ω,\ ‖Z\ ω‖\ \textasciicircum{}\ 2\ ∂Q}
\end{leancode}
\leansource{\texttt{evaluation\_\allowbreak{}bcf\_\allowbreak{}coe} \quad WassersteinMixing/\allowbreak{}SemigroupContinuity.lean, lines 63–69}
\begin{leancode}
\leanline{0}{private\ theorem\ evaluation\_bcf\_coe\ (X\ :\ Lp\ E\ 2\ Q)\ (hX\ :\ Q.map\ X\ =\ μ)}
\leanline{4}{(φ\ :\ E\ →ᵇ\ F)\ :}
\leanline{4}{evaluation\ ({\leanfallback 𝕜}\ :=\ {\leanfallback 𝕜})\ μ\ Q\ X\ hX\ (BoundedContinuousFunction.toLp\ 2\ μ\ {\leanfallback 𝕜}\ φ)}
\leanline{6}{=ᵐ[Q]\ fun\ ω\ ↦\ φ\ (X\ ω)}
\end{leancode}
\leansource{\texttt{evaluation\_\allowbreak{}bcf\_\allowbreak{}norm\_\allowbreak{}sq} \quad WassersteinMixing/\allowbreak{}SemigroupContinuity.lean, lines 72–84}
\begin{leancode}
\leanline{0}{private\ theorem\ evaluation\_bcf\_norm\_sq\ (X\ Y\ :\ Lp\ E\ 2\ Q)}
\leanline{4}{(hX\ :\ Q.map\ X\ =\ μ)\ (hY\ :\ Q.map\ Y\ =\ μ)\ (φ\ :\ E\ →ᵇ\ F)\ :}
\leanline{4}{‖evaluation\ ({\leanfallback 𝕜}\ :=\ {\leanfallback 𝕜})\ μ\ Q\ X\ hX\ (BoundedContinuousFunction.toLp\ 2\ μ\ {\leanfallback 𝕜}\ φ)\ -}
\leanline{6}{evaluation\ ({\leanfallback 𝕜}\ :=\ {\leanfallback 𝕜})\ μ\ Q\ Y\ hY\ (BoundedContinuousFunction.toLp\ 2\ μ\ {\leanfallback 𝕜}\ φ)‖\ \textasciicircum{}\ 2\ =}
\leanline{6}{∫\ ω,\ ‖φ\ (X\ ω)\ -\ φ\ (Y\ ω)‖\ \textasciicircum{}\ 2\ ∂Q}
\end{leancode}
\leanlabel{How the Lean proof follows the argument}
\texttt{evaluation} is actual measure-preserving composition. \texttt{lp\_\allowbreak{}norm\_\allowbreak{}sq\_\allowbreak{}eq\_\allowbreak{}integral} takes the real part of the integrable inner-product formula. The bounded-continuous representative lemmas identify the norm of the actual evaluation difference with the displayed integral.
\end{leanbox}

\section{\texorpdfstring{Extending convergence from bounded continuous functions to all observables}{Extending convergence from bounded continuous functions to all observables}}


\textbf{Claim.} If \(X_i\to X\) in measure and all states have law \(\mu\), then \(\phi(X_i)\to\phi(X)\) in \(L^2\) for bounded continuous \(\phi\), and the same is true for every square-integrable observable.



\textbf{Proof.} Composition with a continuous bounded function preserves convergence in measure. The squared difference has a uniform constant bound, so dominated convergence in measure gives its integral tending to zero. For arbitrary \(f\), approximate it in \(L^2(\mu)\) by a bounded continuous function. All evaluation maps have norm one, so both approximation errors have the same uniform bound; the middle bounded-continuous evaluation difference tends to zero. Norm convergence of states implies convergence in measure, giving the final specializations.

\begin{leanbox}{Lean: \texttt{evaluation\_\allowbreak{}bcf\_\allowbreak{}tendsto\_\allowbreak{}inMeasure}, \texttt{evaluation\_\allowbreak{}tendsto\_\allowbreak{}inMeasure}, \texttt{evaluation\_\allowbreak{}bcf\_\allowbreak{}tendsto} and 1 more}
\leansource{\texttt{evaluation\_\allowbreak{}bcf\_\allowbreak{}tendsto\_\allowbreak{}inMeasure} \quad WassersteinMixing/\allowbreak{}SemigroupContinuity.lean, lines 89–127}
\begin{leancode}
\leanline{0}{theorem\ evaluation\_bcf\_tendsto\_inMeasure\ \{I\ :\ Type*\}\ \{l\ :\ Filter\ I\}\ [l.IsCountablyGenerated]}
\leanline{4}{(X\ :\ I\ →\ Lp\ E\ 2\ Q)\ (Y\ :\ Lp\ E\ 2\ Q)\ (hX\ :\ ∀\ i,\ Q.map\ (X\ i)\ =\ μ)}
\leanline{4}{(hY\ :\ Q.map\ Y\ =\ μ)}
\leanline{4}{(hlim\ :\ TendstoInMeasure\ Q\ (fun\ i\ ω\ ↦\ X\ i\ ω)\ l\ Y)\ (φ\ :\ E\ →ᵇ\ F)\ :}
\leanline{4}{Tendsto\ (fun\ i\ ↦\ evaluation\ ({\leanfallback 𝕜}\ :=\ {\leanfallback 𝕜})\ μ\ Q\ (X\ i)\ (hX\ i)}
\leanline{6}{(BoundedContinuousFunction.toLp\ 2\ μ\ {\leanfallback 𝕜}\ φ))\ l}
\leanline{6}{({\leanfallback 𝓝}\ (evaluation\ ({\leanfallback 𝕜}\ :=\ {\leanfallback 𝕜})\ μ\ Q\ Y\ hY\ (BoundedContinuousFunction.toLp\ 2\ μ\ {\leanfallback 𝕜}\ φ)))}
\end{leancode}
\leansource{\texttt{evaluation\_\allowbreak{}tendsto\_\allowbreak{}inMeasure} \quad WassersteinMixing/\allowbreak{}SemigroupContinuity.lean, lines 132–162}
\begin{leancode}
\leanline{0}{theorem\ evaluation\_tendsto\_inMeasure\ [NormedSpace\ ℝ\ F]\ [μ.WeaklyRegular]\ \{I\ :\ Type*\}\ \{l\ :\ Filter\ I\}}
\leanline{4}{[l.IsCountablyGenerated]\ (X\ :\ I\ →\ Lp\ E\ 2\ Q)\ (Y\ :\ Lp\ E\ 2\ Q)}
\leanline{4}{(hX\ :\ ∀\ i,\ Q.map\ (X\ i)\ =\ μ)\ (hY\ :\ Q.map\ Y\ =\ μ)}
\leanline{4}{(hlim\ :\ TendstoInMeasure\ Q\ (fun\ i\ ω\ ↦\ X\ i\ ω)\ l\ Y)\ (f\ :\ Lp\ F\ 2\ μ)\ :}
\leanline{4}{Tendsto\ (fun\ i\ ↦\ evaluation\ ({\leanfallback 𝕜}\ :=\ {\leanfallback 𝕜})\ μ\ Q\ (X\ i)\ (hX\ i)\ f)\ l}
\leanline{6}{({\leanfallback 𝓝}\ (evaluation\ ({\leanfallback 𝕜}\ :=\ {\leanfallback 𝕜})\ μ\ Q\ Y\ hY\ f))}
\end{leancode}
\leansource{\texttt{evaluation\_\allowbreak{}bcf\_\allowbreak{}tendsto} \quad WassersteinMixing/\allowbreak{}SemigroupContinuity.lean, lines 166–172}
\begin{leancode}
\leanline{0}{theorem\ evaluation\_bcf\_tendsto\ \{I\ :\ Type*\}\ \{l\ :\ Filter\ I\}\ [l.IsCountablyGenerated]}
\leanline{4}{(X\ :\ I\ →\ Lp\ E\ 2\ Q)\ (Y\ :\ Lp\ E\ 2\ Q)\ (hX\ :\ ∀\ i,\ Q.map\ (X\ i)\ =\ μ)}
\leanline{4}{(hY\ :\ Q.map\ Y\ =\ μ)\ (hlim\ :\ Tendsto\ X\ l\ ({\leanfallback 𝓝}\ Y))\ (φ\ :\ E\ →ᵇ\ F)\ :}
\leanline{4}{Tendsto\ (fun\ i\ ↦\ evaluation\ ({\leanfallback 𝕜}\ :=\ {\leanfallback 𝕜})\ μ\ Q\ (X\ i)\ (hX\ i)}
\leanline{6}{(BoundedContinuousFunction.toLp\ 2\ μ\ {\leanfallback 𝕜}\ φ))\ l}
\leanline{6}{({\leanfallback 𝓝}\ (evaluation\ ({\leanfallback 𝕜}\ :=\ {\leanfallback 𝕜})\ μ\ Q\ Y\ hY\ (BoundedContinuousFunction.toLp\ 2\ μ\ {\leanfallback 𝕜}\ φ)))}
\end{leancode}
\leansource{\texttt{evaluation\_\allowbreak{}tendsto} \quad WassersteinMixing/\allowbreak{}SemigroupContinuity.lean, lines 176–182}
\begin{leancode}
\leanline{0}{theorem\ evaluation\_tendsto\ [NormedSpace\ ℝ\ F]\ [μ.WeaklyRegular]\ \{I\ :\ Type*\}\ \{l\ :\ Filter\ I\}}
\leanline{4}{[l.IsCountablyGenerated]\ (X\ :\ I\ →\ Lp\ E\ 2\ Q)\ (Y\ :\ Lp\ E\ 2\ Q)}
\leanline{4}{(hX\ :\ ∀\ i,\ Q.map\ (X\ i)\ =\ μ)\ (hY\ :\ Q.map\ Y\ =\ μ)}
\leanline{4}{(hlim\ :\ Tendsto\ X\ l\ ({\leanfallback 𝓝}\ Y))\ (f\ :\ Lp\ F\ 2\ μ)\ :}
\leanline{4}{Tendsto\ (fun\ i\ ↦\ evaluation\ ({\leanfallback 𝕜}\ :=\ {\leanfallback 𝕜})\ μ\ Q\ (X\ i)\ (hX\ i)\ f)\ l}
\leanline{6}{({\leanfallback 𝓝}\ (evaluation\ ({\leanfallback 𝕜}\ :=\ {\leanfallback 𝕜})\ μ\ Q\ Y\ hY\ f))}
\end{leancode}
\leanlabel{How the Lean proof follows the argument}
The bounded-continuous theorem uses its exact square-norm integral and bounded convergence in measure. \texttt{evaluation\_\allowbreak{}tendsto\_\allowbreak{}in\allowbreak{}Measure} uses density and the common isometry bound. The last two theorems convert norm convergence of states to convergence in measure before applying these results.
\end{leanbox}

\end{document}
